\subsection{A returning syzygy as a graph obstruction}

\begin{lemma}[Proved: a stable returning syzygy on a two-term complex]
\label{CP2:A-stable-syzygy-two-term-amplitude}
Let $A$ be a finite-dimensional rational algebra and $V$ a
finite-dimensional left $A$ module.  Suppose that, for some
$d\ge1$ and $r\ge1$, there is an exact sequence
\[
 0\longrightarrow V^{\oplus r}\oplus P
 \longrightarrow P_{d-1}\longrightarrow\cdots
 \longrightarrow P_0\longrightarrow V\longrightarrow0,
\]
where $P$ and all $P_i$ are projective left $A$ modules.
Let $D$ be a finite-dimensional right $A$ complex in degrees
zero and one.  If $D\otimes_A^{\mathbf L}V$ has bounded
homology, then
\[
 H_j(D\otimes_A^{\mathbf L}V)=0\qquad(j\ge2).
\]
\end{lemma}
\begin{proof}
Put $K_0=V$, let $K_i$ be the successive kernels in the
displayed exact sequence, and identify
$K_d=V^{\oplus r}\oplus P$.  For each $0\le i<d$ there
is a short exact sequence
\[
 0\longrightarrow K_{i+1}\longrightarrow P_i
                  \longrightarrow K_i\longrightarrow0.
\]
Since $P_i$ is projective, $D\otimes_A^{\mathbf L}P_i$
is represented by the ordinary two-term complex
$D\otimes_A P_i$.  Its homology vanishes in degrees at
least two.  The derived-tensor long exact sequence therefore
gives, for every $j\ge2$,
\[
 H_{j+1}(D\otimes_A^{\mathbf L}K_i)
       \cong H_j(D\otimes_A^{\mathbf L}K_{i+1}).
\]
Iterating, with all intermediate indices still at least two,
and then using the projectivity of $P$, gives
\[
 H_{j+d}(D\otimes_A^{\mathbf L}V)
   \cong H_j(D\otimes_A^{\mathbf L}K_d)
   \cong H_j(D\otimes_A^{\mathbf L}V)^{\oplus r}
                 \qquad(j\ge2).
\]
If one group on the right were nonzero, repeated use of this
isomorphism would produce nonzero groups in arbitrarily high
degrees.  Boundedness excludes this, proving the result.
The hypothesis returns to $V$ itself up to projective summands;
no assertion about arbitrary periodic but different syzygies is
being used.
\end{proof}

\begin{theorem}[Proved: a stable-syzygy row character over a graph]
\label{CP2:A-stable-syzygy-row-graph-obstruction}
Let $M$ have a $2\times2$ rectangular minimal band $K$,
put $A=\mathbb QM/\mathbb QK$, and let $V$ be its
one-dimensional left row-augmentation module.  Suppose $V$
admits the exact stable-return sequence of the preceding lemma.
Let a finite nonempty left-$N$, right-$M$ biset have the
equivariant graph comparison and rationally acyclic opposite
corner required in
Theorem~\ref{CP2:A-projective-local-row-graph-obstruction}.
If its collage has bounded rational homology, all rational
degree-two products vanish.

In particular this conclusion holds for $N=R_3$ and every
finite nonempty connecting biset, without any state-count
or second-Betti-number restriction.
\end{theorem}
\begin{proof}
Use the natural augmented graph comparison and collapse the
rationally acyclic corner as in the cited graph-localization
proof.  If the two minimal column-image graphs meet, their
common vertex supplies the bridge augmentation section.
The relative suspension argument already kills the products.
In the disjoint case that proof constructs a two-term right
$A$ complex $D$ and establishes
\[
 H_j(BC;\mathbb Q)=H_{j-2}(D\otimes_A^{\mathbf L}V)
                         \qquad(j\ge4).
\]
The construction and this equality precede the algebraic
localness/projectivity case split in the cited proof.  They
therefore require no such extra hypothesis here.  Boundedness
of the collage makes the derived tensor bounded.  The preceding
lemma kills its degree-two group, so $H_4(BC;\mathbb Q)=0$,
which kills all degree-two products.
The natural graph comparison for a left $R_3$ action is exactly
the one in Corollary~\ref{CP2:A-baseline-right-zero-three-six}.
\end{proof}

\begin{remark*}[The literal first-syzygy certificate]
For a one-dimensional character $\chi:A\to\mathbb Q$, the
canonical free cover $A\to V$ has kernel $J=\ker\chi$.
A sufficient $d=1$ witness consists of idempotents $e_1,\ldots,e_t$
in $A$, an integer $r\ge1$, and a literal $A$-module isomorphism
\[
 V^{\oplus r}\oplus\bigoplus_{i=1}^t Ae_i
                  \xrightarrow{\cong}J.
\]
Checking the algebra table, character, each idempotent, bases
of the principal ideals, every intertwining equation and the
nonzero determinant certifies the full hypothesis.  A search
which fails to find such a witness makes no exclusion.  Longer
return sequences require their own explicit exact module data.
\end{remark*}
